% GENERATED FILE. Edit canonical lessons, then run npm run generate:books.
% canonical-source-hash: fnv1a64:c3969dce6e58b096
% canonical-lessons: MR-C115-hodi, MR-C115-saykal, MR-C115-bajar, MR-C115-gav, MR-C115-kapde

\chapter{Boat, Bicycle, Market, Village, Clothes}
\label{ch:mr-boat-bicycle-market-village-clothes}
\hlchaptermodality{115}

\begin{chapteropening}
\textbf{By the end of this chapter:} \emph{I can say a boat, a bicycle, a market, a village and clothes.}

Complete the last lesson of chapter 115: I can say a boat, a bicycle, a market, a village and clothes.
\end{chapteropening}

\section[\texorpdfstring{\mr{होडी} (hoḍī)}{होडी (hoḍī)}]{\mr{होडी} (hoḍī) — a boat}

\label{lesson:MR-C115-hodi}



\begin{quote}
Before the new one: say the Marathi for an office, then the Marathi for the police.
\end{quote}

\subsection*{You'll want to know: \mr{होडी}}
\textbf{\mr{होडी}} — \emph{hoḍī} — \textquotedblleft{}a boat\textquotedblright{}.

\begin{cousinweb}
\textbf{\mr{होडी}} (\emph{hoḍī}) is a tadbhava of Sanskrit \emph{hoḍa}, \textquotedblleft{}raft, boat\textquotedblright{}, a word from the old word-lists (CDIAL 14174). CDIAL takes it from Dravidian, comparing Kannada \emph{ōḍa}.
\end{cousinweb}

\begin{grammarlens}[title={What you've built}]
One more word for this chapter.
\end{grammarlens}

\subsection*{Your turn}
\begin{itemize}
\raggedright
  \item \emph{Say it:} \emph{hoḍī}
  \item \emph{Say it:} \emph{hoḍī}, once more
  \item \emph{Say it:} say \emph{hoḍī}
  \item \emph{Recall:} say the Marathi for an office, then the Marathi for the police, then say \emph{hoḍī} again
\end{itemize}

\begin{tcolorbox}[breakable,colback=teal!4,colframe=teal!35!black,title={Before you move on}]
What does \mr{होडी} mean? (\textquotedblleft{}A boat\textquotedblright{}.) Say it once more.
\end{tcolorbox}
\section[\texorpdfstring{\mr{सायकल} (sāykal)}{सायकल (sāykal)}]{\mr{सायकल} (sāykal) — a bicycle}

\label{lesson:MR-C115-saykal}



\begin{quote}
Before the new one: say the Marathi for the police, then the Marathi for a boat.
\end{quote}

\subsection*{You'll want to know: \mr{सायकल}}
\textbf{\mr{सायकल}} — \emph{sāykal} — \textquotedblleft{}a bicycle\textquotedblright{}.

\begin{cousinweb}
\textbf{\mr{सायकल}} (\emph{sāykal}) is a loan from English \emph{cycle}.
\end{cousinweb}

\begin{grammarlens}[title={What you've built}]
One more word for this chapter.
\end{grammarlens}

\subsection*{Your turn}
\begin{itemize}
\raggedright
  \item \emph{Say it:} \emph{sāykal}
  \item \emph{Say it:} \emph{sāykal}, once more
  \item \emph{Say it:} say \emph{sāykal}
  \item \emph{Recall:} say the Marathi for the police, then the Marathi for a boat, then say \emph{sāykal} again
\end{itemize}

\begin{tcolorbox}[breakable,colback=teal!4,colframe=teal!35!black,title={Before you move on}]
What does \mr{सायकल} mean? (\textquotedblleft{}A bicycle\textquotedblright{}.) Say it once more.
\end{tcolorbox}
\section[\texorpdfstring{\mr{बाजार} (bājār)}{बाजार (bājār)}]{\mr{बाजार} (bājār) — a market}

\label{lesson:MR-C115-bajar}



\begin{quote}
Before the new one: say the Marathi for a boat, then the Marathi for a bicycle.
\end{quote}

\subsection*{You'll want to know: \mr{बाजार}}
\textbf{\mr{बाजार}} — \emph{bājār} — \textquotedblleft{}a market\textquotedblright{}.

\begin{cousinweb}
\textbf{\mr{बाजार}} (\emph{bājār}) is a loan, usually traced to Persian \emph{bāzār}, \textquotedblleft{}market\textquotedblright{}. Marathi writes the \emph{z} as \textbf{\mr{ज}}.
\end{cousinweb}

\begin{grammarlens}[title={What you've built}]
One more word for this chapter.
\end{grammarlens}

\subsection*{Your turn}
\begin{itemize}
\raggedright
  \item \emph{Say it:} \emph{bājār}
  \item \emph{Say it:} \emph{bājār}, once more
  \item \emph{Say it:} say \emph{bājār}
  \item \emph{Recall:} say the Marathi for a boat, then the Marathi for a bicycle, then say \emph{bājār} again
\end{itemize}

\begin{tcolorbox}[breakable,colback=teal!4,colframe=teal!35!black,title={Before you move on}]
What does \mr{बाजार} mean? (\textquotedblleft{}A market\textquotedblright{}.) Say it once more.
\end{tcolorbox}
\section[\texorpdfstring{\mr{गाव} (gāv)}{गाव (gāv)}]{\mr{गाव} (gāv) — a village}

\label{lesson:MR-C115-gav}



\begin{quote}
Before the new one: say the Marathi for a bicycle, then the Marathi for a market.
\end{quote}

\subsection*{You'll want to know: \mr{गाव}}
\textbf{\mr{गाव}} — \emph{gāv} — \textquotedblleft{}a village\textquotedblright{}.

\begin{cousinweb}
\textbf{\mr{गाव}} (\emph{gāv}) is a tadbhava, a word inherited from Sanskrit and reshaped by sound change. It comes from \emph{grāma}, \textquotedblleft{}troop, village\textquotedblright{} (CDIAL 4368).
\end{cousinweb}

\begin{grammarlens}[title={What you've built}]
One more word for this chapter.
\end{grammarlens}

\subsection*{Your turn}
\begin{itemize}
\raggedright
  \item \emph{Say it:} \emph{gāv}
  \item \emph{Say it:} \emph{gāv}, once more
  \item \emph{Say it:} say \emph{gāv}
  \item \emph{Recall:} say the Marathi for a bicycle, then the Marathi for a market, then say \emph{gāv} again
\end{itemize}

\begin{tcolorbox}[breakable,colback=teal!4,colframe=teal!35!black,title={Before you move on}]
What does \mr{गाव} mean? (\textquotedblleft{}A village\textquotedblright{}.) Say it once more.
\end{tcolorbox}
\section[\texorpdfstring{\mr{कपडे} (kapḍe)}{कपडे (kapḍe)}]{\mr{कपडे} (kapḍe) — clothes}

\label{lesson:MR-C115-kapde}



\begin{quote}
Before the new one: say the Marathi for a market, then the Marathi for a village.
\end{quote}

\subsection*{You'll want to know: \mr{कपडे}}
\textbf{\mr{कपडे}} — \emph{kapḍe} — \textquotedblleft{}clothes\textquotedblright{}.

\begin{cousinweb}
\textbf{\mr{कपडे}} (\emph{kapḍe}) is the plural of \emph{kapḍā}, \textquotedblleft{}cloth\textquotedblright{}, which goes back, like Marathi \emph{kāpaḍ}, \textquotedblleft{}cloth\textquotedblright{}, to Sanskrit \emph{karpaṭa}, \textquotedblleft{}patched garment, rag\textquotedblright{} (CDIAL 2871). CDIAL reports Przyluski's view that the Sanskrit word comes from Austroasiatic.
\end{cousinweb}

\begin{grammarlens}[title={What you've built}]
One more word for this chapter.
\end{grammarlens}

\subsection*{Your turn}
\begin{itemize}
\raggedright
  \item \emph{Say it:} \emph{kapḍe}
  \item \emph{Say it:} \emph{kapḍe}, once more
  \item \emph{Say it:} say \emph{kapḍe}
  \item \emph{Recall:} say the Marathi for a market, then the Marathi for a village, then say \emph{kapḍe} again
\end{itemize}

\begin{tcolorbox}[breakable,colback=teal!4,colframe=teal!35!black,title={Before you move on}]
What does \mr{कपडे} mean? (\textquotedblleft{}Clothes\textquotedblright{}.) Say it once more.
\end{tcolorbox}
